\section{Introduction}
\label{sec:intro}

A satellite in low Earth orbit (LEO) can observe far more of the ground
than it can ever report. A single imaging pass over a wildfire, a flood,
or an agricultural region can generate far more raw pixels than the
satellite's downlink budget can move to a ground station before the next
useful decision has to be made — and a LEO satellite is only in radio
contact with any one ground station for a few minutes at a time before it
drops below the horizon. Someone, or something, has to decide which
images matter most, how aggressively to compress each one, and when to
send it. \emph{Orbital Edge Computing} (OEC) is the idea of making some of
that decision on board the satellite itself — compressing and
prioritizing at the edge of the network, rather than shipping raw data
down and sorting it out on the ground.

This paper studies that decision problem end to end: an on-board
residual-quantized variational autoencoder (RQ-VAE) compresses each image
at one of several selectable \emph{depths} (more refinement passes, larger
payload, higher fidelity), and a rolling-horizon Model Predictive Control
(MPC) scheduler decides, for every task competing for downlink capacity,
which depth to use and when and along which route to transmit it, under a
constellation topology and ground-station contact schedule that changes
every scheduling slot.

\paragraph{The problem with grading a scheduler on how the image looks.}
The natural way to score a compression-and-scheduling policy is to ask
whether the reconstructed image still \emph{looks} like the original —
standard fidelity metrics such as PSNR and LPIPS answer exactly that
question. But for a mission whose actual goal is to extract information
from the imagery (what land-cover class is this, where is the flood
boundary), \emph{looking similar} and \emph{remaining useful for the task}
are not the same property, and this paper's first and most consequential
finding is that they diverge sharply here: pixel-fidelity quality changes
by only $12\%$ across the compression-depth range we test, while the
imagery's usefulness for a real downstream semantic-segmentation task
changes by $247\%$ over the identical range (\S\ref{sec:downstream}). A
scheduler rewarded on the flat proxy has almost nothing to gain by
reasoning carefully about depth; a scheduler rewarded on the true
task-relevant signal has a great deal to gain. This single substitution —
replacing a pixel-fidelity proxy with a measured downstream-task quality
signal — changes MPC's measured advantage over the best fixed-depth
baseline from a marginal $1.2\%$ to a substantial $32.2\%$ in an otherwise
identical experiment (\S\ref{sec:utility-results}).

\paragraph{Contributions.} Concretely, this paper:
\begin{enumerate}
\item grounds a scheduling utility function's quality term in
  \emph{measured} downstream semantic-segmentation mIoU rather than pixel
  fidelity, and shows this changes not just the numbers but the
  qualitative conclusion about whether intelligent scheduling matters
  (\S\ref{sec:compression}, \S\ref{sec:utility-results});
\item presents a single \emph{unified utility function} — quality,
  tardiness, resource cost, and a fairness floor, combined so that the
  resulting optimization stays a mixed-integer \emph{linear} program with
  no auxiliary binary variables — that replaces five previously
  disagreeing scoring definitions used across the scheduler, the offline
  bound, and two prior hierarchical variants (\S\ref{sec:formulation});
\item evaluates ten schedulers — four fixed-depth and one adaptive greedy
  baseline, a flat rolling-horizon MPC, a congestion-aware routing
  variant, and two ways of coupling a dedicated routing optimizer to the
  depth decision (hierarchical and two-level/peer) — against a
  three-tiered offline optimality bound, under two link-capacity regimes
  designed to make different resources binding (\S\ref{sec:setup},
  \S\ref{sec:utility-results});
\item shows, from the resulting depth-selection decisions and a
  link-rate sensitivity sweep, that \emph{admission control} — deciding
  early to drop a task rather than let it starve — is what separates the
  strongest and weakest MPC variants, not how cleverly the kept tasks are
  routed (\S\ref{sec:decisions}, \S\ref{sec:sweeps}); and
\item reports a structural, not tunable, negative result: freezing a
  concrete route for a multi-minute planning epoch is usable at its own
  executed slot only a small fraction of the time, because ground-station
  contact windows in this constellation are shorter than any freeze period
  worth having (\S\ref{sec:formulation}, \S\ref{sec:utility-results}); and
  we additionally evaluate a learned (PPO) scheduling policy as a
  compute-cost point of comparison to MPC (\S\ref{sec:rl}).
\end{enumerate}

\paragraph{Roadmap.} \S\ref{sec:related-work} situates this work relative
to orbital edge computing, learned image compression, and MPC-based
network scheduling. \S\ref{sec:formulation} gives the system model and the
full scheduling formulation, including every MPC variant evaluated later.
\S\ref{sec:setup} specifies the simulated constellation, link-capacity
regimes, and task workload used to reproduce every subsequent result.
\S\ref{sec:compression}--\S\ref{sec:rl} present results: compression and
downstream-task quality, the unified-utility scheduler comparison,
depth-selection decision distributions, link-rate sensitivity, and the
learned-policy baseline, in that order. \S\ref{sec:conclusion} summarizes
and discusses what remains open.
